\documentclass[11pt]{article}
\usepackage{amsmath,amssymb,amsthm,geometry,hyperref}
\geometry{margin=1in}
\newtheorem{lemma}{Lemma}
\newtheorem{proposition}[lemma]{Proposition}
\newtheorem{corollary}[lemma]{Corollary}
\theoremstyle{remark}\newtheorem{remark}[lemma]{Remark}
\newcommand{\R}{\mathcal R}
\newcommand{\DeltaT}{\Delta(T)}
\title{Tree gluing and an analytic replacement for pentagon factors}
\author{Research candidate; main manuscript unchanged}
\date{10 October 2026}
\begin{document}
\maketitle

All graphs below are finite and simple. A path factor is a spanning
subgraph whose components are paths; its paths need not be induced.
This candidate preserves the assertions and palette thresholds of the
main manuscript. It replaces part of the proof infrastructure, not the
protected-edge identities or the bad-block lifting lemma.

\section*{1. The common gluing lemma}

\begin{lemma}[Reserved-edge tree gluing]\label{lem:tree-gluing}
Let a graph $X$ have a vertex partition $(W_v)_{v\in V(T)}$, where
$T$ is a tree. For each $v$, let $C_v$ be a Hamilton cycle on $W_v$.
For every edge $vw\in E(T)$ choose edges
$x_{vw}^vy_{vw}^v\in E(C_v)$ and
$x_{vw}^wy_{vw}^w\in E(C_w)$ such that
$x_{vw}^vx_{vw}^w$ and $y_{vw}^vy_{vw}^w$ are edges of $X$.
Suppose the edges chosen at each fixed $v$ are distinct.
Then the graph obtained from the disjoint union of the $C_v$ by
these two-edge switches is a Hamilton cycle of $X$.
Every edge of $C_v$ not selected for a switch is retained.
\end{lemma}
\begin{proof}
Root $T$ and add its vertices in parent-before-child order. Initially
take the root cycle. At every subsequent step the parent belongs to
the one cycle already constructed and the child still has its original
cycle. The selected parent edge is still present, because no earlier
switch selected that edge; the child edge is also present. Deleting
these two edges leaves two paths, and the two specified cross edges
join them to one cycle. Induction covers every block once. For a
one-vertex tree no switch is needed. Only selected edges are removed.
\end{proof}

\begin{remark}
Selected edges in one reference cycle may share an endpoint. They
must be distinct as edges, not vertex-disjoint. Counting distinct
deleted reference edges remains valid after all earlier switches:
the reference edges still present are exactly those not yet selected.
New cross edges are never required as later ports.
\end{remark}

\begin{corollary}[Fresh-child port counting]\label{cor:port-counting}
Root $T$. Suppose each tree interface permits at least $\DeltaT$
different edges of the parent reference cycle, each with a compatible
edge in the fresh child reference cycle. Then the edges required in
Lemma~\ref{lem:tree-gluing} can be chosen greedily.
More generally, partition the possible reference edges in each block
into disjoint row classes. If $b_{v,z}$ incident tree interfaces use
row $z$, it suffices that each such interface offers at least
$b_{v,z}$ distinct edges in that parent row.
\end{corollary}
\begin{proof}
Before a particular parent--child interface is processed, at most
$\deg_T(v)-1$ reference edges have been selected in its parent block.
In the row version at most $b_{v,z}-1$ have been selected in that row.
Choose a different offered edge. The child is fresh, so its compatible
edge has not been selected. Continue recursively. This includes the
child's incoming edge in the count when it later acts as a parent.
\end{proof}

\begin{remark}[Cycles realized after port selection]
The tree lemma also applies if the cycles are constructed only after
the ports have been selected. It is enough to prove, for each block,
that its complete selected port set lies on a Hamilton cycle of that
block. This is the form needed for the $(2,3,2)$ bad blocks. It must
not be replaced by a fixed-reference-cycle counting assertion unless
the requisite cycle has already been constructed.
\end{remark}

\section*{2. A single analytic input replaces listed path factors}

\begin{lemma}[Cut bound for complete list recolouring]\label{lem:cut}
Let $H$ be triangle-free, with fixed finite colour lists at its
vertices. Let $Y$ be a connected component of its complete list
recolouring graph, retaining all legal single-vertex moves. For
every nonempty $S\subseteq V(Y)$,
\[
                         c(Y-S)\le2|S|.
\]
\end{lemma}
\begin{proof}
The case $S=V(Y)$ is immediate. For a component $D$ of $Y-S$, put
$N_S(D)=\{s\in S:N_Y(s)\cap V(D)\ne\varnothing\}$ and
$b(D)=|N_S(D)|$. Connectivity makes $b(D)>0$.
For each $s\in N_S(D)$ choose a representative
$i_{s,D}\in N_Y(s)\cap V(D)$. Representatives chosen for different
components at a fixed $s$ are nonadjacent, so their moves change
different original vertices.

Form a graph $M_s$ on the components incident with $s$, joining two
when their representative moves commute. Two legal moves at distinct
vertices fail to commute precisely when these original vertices are
adjacent and their new colours agree. Thus an independent set of
$M_s$ would give a clique of $H$, and $\alpha(M_s)\le2$.
For a commuting pair $D,E$, the double-change state is adjacent to
both representatives. It belongs to $Y$; it must lie in $S$, since
otherwise it joins $D$ to $E$ in $Y-S$. For fixed $s,D$, the different
partners $E$ give distinct such states in $N_S(D)\setminus\{s\}$.
Consequently $d_{M_s}(D)+1\le b(D)$.

For every finite graph $M$,
$\sum_{u\in V(M)}1/(d_M(u)+1)\le\alpha(M)$: in a uniform random
ordering select each vertex preceding all its neighbours; this is an
independent set, with the stated expected size. Hence
\[
 \sum_{D:s\in N_S(D)}\frac1{b(D)}
 \le\sum_{D\in V(M_s)}\frac1{d_{M_s}(D)+1}\le2.
\]
Sum over $s\in S$. Every component $D$ contributes $b(D)$ copies of
$1/b(D)$, giving exactly $c(Y-S)$.
\end{proof}

\begin{lemma}[$P_2/P_3$ factors]\label{lem:path-factor}
Every nontrivial component $Y$ in Lemma~\ref{lem:cut} has a
$\{P_2,P_3\}$ factor.
\end{lemma}
\begin{proof}
For any nonempty independent set $I$ of $Y$, its neighbourhood is
nonempty, and deleting it isolates all vertices of $I$.
Lemma~\ref{lem:cut} gives $|I|\le2|N_Y(I)|$.
For arbitrary $A\subseteq V(Y)$, write
$\Gamma(A)=\bigcup_{a\in A}N_Y(a)$, including neighbours lying
inside $A$, and put $I=A\setminus\Gamma(A)$. Then $I$ is independent
and $\Gamma(I)\subseteq\Gamma(A)\setminus A$. Thus
\[
 |A|=|A\cap\Gamma(A)|+|I|
 \le |A\cap\Gamma(A)|+2|\Gamma(A)\setminus A|
 \le2|\Gamma(A)|.
\]
Apply Hall's theorem to a bipartite graph having one left copy and
two right copies of each vertex, with adjacency inherited from $Y$.
There is a map $f:V(Y)\to V(Y)$ with $vf(v)\in E(Y)$ and at most
two preimages at each vertex. The digraph $v\to f(v)$ has outdegree
one, indegree at most two, no loops, and each weak component consists
of a directed cycle of length at least two with incoming trees.

While a noncycle vertex has children, choose one furthest from its
cycle. Its one or two children are leaves. Delete this vertex with
all its children as a $P_2$ or $P_3$. No surviving vertex points to a
deleted vertex, so repeat. Afterwards noncycle vertices are leaves
directly attached to cycle vertices, with at most one leaf per cycle
vertex. If there are attached leaves, split the directed cycle into
segments starting at a leaf-bearing vertex $z$ and ending immediately
before the next such vertex. If the $r$ subsequent leaf-free cycle
vertices in a segment are even in number, pair $z$ with its leaf and
pair the $r$ vertices consecutively. If $r$ is odd, take the leaf,
$z$, and the next cycle vertex as a $P_3$, then pair the remainder.
These segments partition the cycle and its leaves. If there are no
leaves, partition an even cycle into edges; for an odd cycle take one
three-vertex path and pair the remainder. A directed two-cycle is an
edge and obeys the same rules. Together with the deleted paths this
gives the required factor.
\end{proof}

\begin{proposition}[Physical connection load three]\label{prop:load}
For any prescribed path factor $F$ of a nontrivial $Y$ in
Lemma~\ref{lem:cut}, whose paths all have order at least two, there
are $c(F)-1$ edges joining these paths in a tree after contraction,
such that every actual vertex of $Y$ is incident with at most three
of these added edges.
\end{proposition}
\begin{proof}
Use Theorem 3.5 of Ellingham--Zha~\cite{EZ}, with $g=2,k=3$.
Its relevant hypothesis is
$c(Y-S)<(k-1)|S|+3=2|S|+3$ for every $S\subseteq V(Y)$.
For nonempty $S$ this follows from Lemma~\ref{lem:cut}; for $S$ empty,
connectivity gives $1<3$. Every component of the prescribed spanning
subgraph $F$ has at least two vertices. The theorem supplies exactly
$c(F)-1$ cross-component edges and bounds their incidence at each
physical vertex by three. Their contraction is connected with one
fewer edges than vertices, hence a tree. This theorem bounds added
edges, not total degree including the path edges.
\end{proof}

\section*{3. The simplified product lemma}

\begin{lemma}[Protected products from two seeds]\label{lem:product}
Let $A$ be a graph, and let $K,P\subseteq E(A)$ be disjoint, with
$|P|=3$. Suppose $A\square P_j$ has a Hamilton cycle retaining every
copy of every edge of $K\cup P$ in every row, for $j=2,3$.
Let $B$ be any nontrivial connected component of a complete fixed-list
recolouring graph of a triangle-free graph. Then $A\square B$ has
a Hamilton cycle retaining a copy of each edge of $K$ in every
physical $B$ row.
\end{lemma}
\begin{proof}
By Lemma~\ref{lem:path-factor}, partition $B$ into $P_2$ and $P_3$
paths and use the corresponding seed cycle on each path product.
Proposition~\ref{prop:load} connects the paths in a tree, with at
most three connecting edges at any physical state $b$ of $B$.
Each connecting edge $bb'$ permits a switch at any of the three
copies of edges in $P$, in rows $b,b'$.
Process the path blocks in parent-before-child order. Before the
next switch, at most two of these three edges have been removed
in the parent row; the entire child block is fresh. Choose an unused
one and use the corresponding child edge. Lemma~\ref{lem:tree-gluing}
then gives one spanning cycle. All copies of $K$ survive, since only
copies of edges of $P$ were selected. Notice that a path block can
have large tree degree; the bound of three applies separately to
each of its two or three physical rows.
\end{proof}

\begin{remark}[Variant with fewer external dependencies]
Without Proposition~\ref{prop:load}, choose any tree in the quotient
of the path factor. If $\Delta(B)\le D$, each physical state has at
least one incident factor edge and can witness at most $D-1$ added
tree edges. The same proof works with $|P|=D-1$. For listed pentagons
the existing elementary $D=7$ estimate therefore retains six ports.
This variant still eliminates the finite factor words and all
arbitrary-length $A\square P_r$ constructions; the three-port
variant also eliminates the need for the degree-seven estimate.
\end{remark}

\begin{corollary}[Listed pentagon application]\label{cor:pentagon}
Keep all protected sets and pools in the current manuscript unchanged.
For every two of its 23 pentagon component types $A,B$, the conclusion
of its protected-product lemma follows from Lemma~\ref{lem:product}.
Neither the eleven printed path factors nor product constructions
for row counts greater than three are required.
\end{corollary}
\begin{proof}
Every proper four-colouring of $C_5$ repeats a colour, since there
are five vertices. Repeated vertices are nonadjacent, hence distance
two on $C_5$; their common neighbour has two equally coloured
neighbours. A list of size at least three leaves at least two colours
available there, so the state has a legal move. Thus every pentagon
component $B$ is nontrivial. It is a component of the complete list
recolouring graph of the triangle-free graph $C_5$, as required.

For a Hamilton component $A$, choose three of the last six pool
edges as $P$, and keep the original $K_A$. Construct $A\square P_2$
by switching the first pool edge between two copies of $C_A$.
Construct $A\square P_3$ by switching the first and second pool
edges at its successive interfaces. The middle row uses two distinct
edges. Lemma~\ref{lem:tree-gluing} gives seed cycles retaining
$K_A\cup P$ in every row. For either hard component, its printed
two- and three-row seeds already retain all of $D_A$ in every row;
again take any three of the last six pool edges and $K=K_A$.
Lemma~\ref{lem:product} now applies.
\end{proof}

\begin{remark}[Optional smaller pools; not applied to the manuscript]
For a Hamilton component one may retain just the first two pool edges
and three of the remaining six: a five-edge pool suffices.
For a hard component a three-edge pool chosen from the last six
suffices; the former seed-block joining edge is unnecessary.
Define the new protected set by removing only the smaller pool from
$E(C_A)$ or $D_A$, respectively. Then $K_A^{\rm old}\subseteq
K_A^{\rm new}$, so every old common-edge witness remains valid.
The seed argument above also retains the enlarged protected sets.
This is a mathematical monotonicity argument, not a claim that an
unchanged certificate reader accepts altered pool formats. Keeping
the original pools avoids that software-format change entirely.
No minimality assertion about three ports is made.
\end{remark}

\section*{4. How the same gluing lemma shortens the other lifts}

\paragraph{Two separated pentagons.}
Keep the existing proof that the bottom component graph is connected,
and root a spanning tree of it. It has maximum degree at most eight:
there are at most four product components in either adjacent layer.
Assign the restoration maps recursively from parent to child, as in
the current manuscript. The protected-edge identities give a common
left edge and the overlap lemma gives at least ten common right
states; hence the two reference product cycles have at least ten
distinct common internal edges. The simplified product lemma supplies
these cycles with all required left edges in all physical right rows.
Before any next parent--child switch at most seven reference edges
of the parent component have been removed. The fresh child has lost
none. Thus Corollary~\ref{cor:port-counting} chooses a surviving
common edge. Its two cross edges are actual bottom recolourings
holding the complete local state fixed. Lemma~\ref{lem:tree-gluing}
finishes the restoration. This is one application of the generic
lemma, not a new proof of component-graph connectivity or compatibility.

\paragraph{The $(2,3,2)$ path.}
Keep the existing good-cycle and safe-common-edge identities, the
bad-block lifting lemma, and the argument partitioning a bottom
Hamilton path into singleton good fibres and bad blocks with their
good endpoints. Select an interface edge at each successive block
boundary. At a singleton block the incoming and outgoing ports must
be different edges of its fixed good cycle; two safe candidates make
this possible. In a bad block, the two ports lie in different end
fibres, so they are distinct physical edges, and the bad-block lemma
realizes a Hamilton cycle containing both. After these cycles have
been constructed, apply Lemma~\ref{lem:tree-gluing} with the block
path as $T$. This covers the case of one block without any switch.
The path lemma cannot be reduced just to fixed-reference port counting:
its bad-block cycles are constructed in response to the reserved ports.

\section*{5. A short endblock structural lemma}

\begin{lemma}[Pentagon in an irreducible endblock]\label{lem:endblock}
Let $G$ be planar of girth at least five, with minimum degree at least
two and every block internally subcubic. Suppose none of the following
occurs: adjacent degree-two vertices; a degree-$(2,3,2)$ path; two
degree-two vertices at distance at least three. Let $B$ be an endblock
with unique cutvertex $x$. Then $B$ has a facial pentagon avoiding
$x$, containing at most one vertex of degree two in $G$.
\end{lemma}
\begin{proof}
The endblock cannot be a bridge, whose noncut endpoint would be a
leaf. Thus $B$ is 2-connected, all its degrees are two or three, and
its noncut vertices have the same degree in $B$ and $G$.
Let $q$ count the degree-two vertices of $B-x$ and put
$\varepsilon=[d_B(x)=2]$. If $q\ge2$, any two such vertices are
nonadjacent and have distance two. A common neighbour of degree two
would give adjacent degree-two vertices, and one of degree three
would give the excluded $(2,3,2)$ path. The external neighbours of
such a path are distinct and outside it by girth five. Thus their
common neighbour has degree at least four in $G$, and can only be
$x$. Every one of the $q$ vertices is consequently adjacent to $x$,
giving $q+\varepsilon\le3$. The same inequality holds when $q\le1$.

Euler's identity and the face-length sum give
\[
 f_5(B)=12-2(q+\varepsilon)+
             \sum_{|F|\ge7}(|F|-6)\ge6.
\]
Faces of length five are simple pentagons, while at most
$d_B(x)\le3$ faces meet $x$. Choose a pentagon avoiding $x$.
If $q\le1$ it has at most one degree-two vertex. If $q\ge2$, every
cycle through one of those vertices uses its edge to $x$, so the
chosen pentagon contains none of them.
\end{proof}

\begin{corollary}[Two endblocks]
If $G$ in Lemma~\ref{lem:endblock} is connected and has more than
one block, it has two disjoint pentagons with no edge between them,
each with at most one degree-two vertex and all other vertices of
degree three in $G$.
\end{corollary}
\begin{proof}
The block-cut tree has at least two endblocks. Apply the lemma in
two different endblocks. Their chosen pentagons contain no cutvertex;
each such vertex belongs to its one block only. Thus the pentagons
are disjoint and have no edge between them.
\end{proof}

\section*{6. Source locations and remaining dependencies}

The source manuscript is
\texttt{planar\_recoloring\_hamilton\_thresholds.tex}. The following
line spans refer to the first source read in this research turn,
before concurrent edits shifted later sections; the named lemma and
certificate tags are the stable locators. They are navigation aids,
not a claim that the live manuscript still has these line numbers:
\begin{itemize}
\item $1471$--$1697$: $(2,3,2)$ fibres, safe edges, bad blocks, and lifting.
\item $1716$--$1743$: finite listed pentagon factor lemma.
\item $1745$--$1808$: protected cycles, seeds, pools, and 313 interfaces.
\item $1810$--$1854$: degree-seven estimate and protected products.
\item $1860$--$1951$: overlap, component graph, and two-pentagon lift.
\item $2095$--$2127$: endblock selection in the block-subcubic proof.
\item $2400$--$2536$: the eleven literal factors. The $30$ words for
the two constant $C_{15}$ components are still needed as their cycle
specifications unless they are restated elsewhere. The other $560$
factor words cease to be required for this product proof.
\end{itemize}

The cut and path-factor arguments above reuse the local research
notes \nolinkurl{trianglefree_list_recolouring_cut_bound_20260915_zh.md}
and \nolinkurl{trianglefree_list_recolouring_path_factor_20260915_zh.md}
under the existing graph-theory project's high-degree notes. They are
included here in full at the relevant mathematical scope rather than
treated as software verdicts. The only additional nonclassical
external theorem used by the three-port version is the author-hosted
Ellingham--Zha theorem cited below. The six-port version does not
need that theorem.

This candidate does not eliminate the finite cycle/seed specifications
or the 313 protected-edge interface identities. Individual pentagon
fibres need not be Hamiltonian: the constant fibre has two components,
and the $00112$ fibre has degree-one states $12303$ and $13203$.
Nor does state overlap alone imply overlap of the chosen cycle edges.
The bad-block lemma and the safe-interface table for the three-vertex
path also remain genuine local inputs.

\begin{thebibliography}{1}
\bibitem{EZ} M. N. Ellingham and X. Zha,
\emph{Toughness, trees and walks}, author-hosted manuscript,
21 July 1998, minor revisions 25 July 1999, Theorem 3.5 and the
definition of a $k$-quasitree in Section 3.
\url{https://math.vanderbilt.edu/ellingmn/paper/ttw/b-eps.pdf}.
The statement and the relevant $g=2,k=3$ branch were read directly.
\end{thebibliography}
\end{document}
